\subsection{部分算子}

在本节中， K 是交换域。

我们回顾（E, II, §3, p. 9–10）：\(\text{graph}^{(2)}\) 是这样一个集合，其元素都是有序偶。设 A 与 B 是集合， A 与 B 之间的一个 \(\text{correspondence}\) 是一个三元组 \((\Gamma, A, B)\) ，其中 \(\Gamma\) 是包含在 A × B 中的一个图像；它的定义集（也称为 \(\text{domain}\) ）是 \(\text{pr}_1(\Gamma)\) ，而它的值集是 \(\text{pr}_2(\Gamma)\) 。一个对应是一个 \(\text{function}\) （E, II, p. 13，定义 9），如果它的图像是泛函的，且它的出发集与它的定义域重合。泛函图像的每个子集都是泛函图像。

定义 1. --- 设 E 与 F 是 K 上的向量空间。从 E 到 F 的部分算子 u 是 E 与 F 之间的一个对应 (Γ, E, F) ，它满足下述条件：

\begin{enumerate}
\def\labelenumi{(\roman{enumi})}
\item
  图像 \(\Gamma\) 是 \(\mathrm{E} \times \mathrm{F}\) 的一个向量子空间；
\item
  图像 \(\Gamma\) 是泛函的。
\end{enumerate}

若 E = F ，就说 u 是 E 上的部分算子。

设 u 是从 E 到 F 的部分算子。对应 u 的图像 \(\Gamma\) 称为部分算子 u 的图像，也记作 \(\Gamma_{u}\) 。我们用 \(\mathcal{P}(\mathrm{E};\mathrm{F})\) 表示从 E 到 F 的部分算子所成的集；我们简记 \(\mathcal{P}(\mathrm{E}) = \mathcal{P}(\mathrm{E};\mathrm{E})\) 。

给出一个从 E 到 F 的部分算子，相当于给出 E 的一个向量子空间 D 与一个从 D 到 F 的线性映射 u ；相应的部分算子是对应 \((\Gamma, \mathrm{E}, \mathrm{F})\) ，其中 \(\Gamma \subset D \times F\) 是 u 的图像。

部分算子 \(u\) 的定义域简称为 \(u\) 的定义域，并记作 \(\operatorname{dom}(u)\) 。

每个从 E 到 F 的线性映射 u 都是从 E 到 F 的部分算子。

若 D ⊂ E 是向量子空间，我们将用 1 \(_{D}\) 表示以 D 为定义域、在 D 上为恒等映射的部分算子，即对应 (Δ \(_{D}\) , E, E) ，其中 Δ \(_{D}\) 是 D × D 的对角线（E, II, p. 13，定义 8）。我们将用 \(0_{\mathrm{D}}\) 表示以 D 为定义域、在 D 上为零的部分算子，即对应 \((\mathrm{D} \times \{0\}, \mathrm{E}, \mathrm{F})\) 。

部分算子 \(u = (\Gamma, \mathrm{E}, \mathrm{F})\) 与 \(u' = (\Gamma', \mathrm{E}, \mathrm{F})\) （都是从 E 到 F 的）相等，当且仅当 \(\operatorname{dom}(u) = \operatorname{dom}(u')\) ，且从 \(u\) 到 F 中的线性映射 \(u'\) 与 \(\operatorname{dom}(u)\) 重合。

依照 E, II, §3 ，定义下列概念：

\begin{enumerate}
\def\labelenumi{(\roman{enumi})}
\tightlist
\item
  u 下 E 的子集 A 的像是 F 的子集 \(u(A \cap D)\) ；简记为 \(u(A)\) 。 u 下 F 的子集 B 的原像是 D 的子集 \(\overline{u}(B)\) 。
\end{enumerate}

若 A （相应地， B ）是 E （相应地， F ）的向量子空间，则它在 u 下的像（相应地，原像）是 F （相应地， E ）的向量子空间。

 \(u\) 的像是 \(u(\mathrm{D})\) 的向量子空间 \(\mathrm{F}\) ，也记作 \(\operatorname{Im}(u)\) 。称 \(u\) 为满射部分算子，如果 \(\operatorname{Im}(u) = \mathrm{F}\) 。 \(u\) 的核是 \(u^{-1}(\{0\})\) 的向量子空间 \(\mathrm{E}\) ，也记作 \(\operatorname{Ker}(u)\) 。 \(u\) 的核化为 0 ，当且仅当从 \(u\) （ \(\operatorname{dom}(u)\) 到 \(\mathrm{F}\) 中的）线性映射是单射的。这时称 \(u\) 是单射的。若 \(u\) 既单射又满射，就称它是双射的。

\begin{enumerate}
\def\labelenumi{(\roman{enumi})}
\setcounter{enumi}{1}
\item
  若 E 、 F 与 G 是 K 上的向量空间，且 \(u = (\Gamma, \mathrm{E}, \mathrm{F})\) ， \(v = (\Gamma', \mathrm{F}, \mathrm{G})\) 分别是从 E 到 F 与从 F 到 G 的部分算子，则复合对应 \(v \circ u = (\Gamma' \circ \Gamma, \mathrm{E}, \mathrm{G})\) 是从 E 到 G 中的部分算子。它的定义域是 \(u^{-1}(\mathrm{dom}(v))\) 。若 H 是 K 上的向量空间，且 \(w = (\Gamma'', \mathrm{G}, \mathrm{H})\) 是从 G 到 H 中的部分算子，则有 \(w \circ (v \circ u) = (w \circ v) \circ u\) 。有时把 \(v \circ u\) 简写作 vu 。
\item
  特别地，对每个从 E 到 F 中的部分算子 \(u\) 以及每个 \(a \in \mathrm{K}\) ，部分算子 \(au = (a1_{\mathrm{F}}) \circ u\) 与 \(ua = u \circ (a1_{\mathrm{E}})\) 有定义。若 \(a \neq 0\) 或者 \(u\) 的定义域等于 E ，则二者相等；我们有 \(u0 = 0_{\mathrm{E}}\) 与 \(0u = 0_{\mathrm{dom}(u)}\) 。
\end{enumerate}

设 E 是向量空间。由前所述，集合 \(\mathcal{P}(\mathrm{E})\) （赋有由 \((u,v)\mapsto u\circ v\) 所定义的复合律）是一个结合的含幺群胚（A, I, p. 4，定义 5 与 A, I, p. 12，定义 2），其单位元为 \(1_{\mathrm{E}}\) 。对每个 \(n\in \mathbf{N}\) ，我们用 \(u^n\) 表示复合 \({}^n u\) （A, I, p. 13）。

此外，我们定义下列概念：

\begin{enumerate}
\def\labelenumi{(\roman{enumi})}
\item
  若 \(u = (\Gamma, \mathrm{E}, \mathrm{F})\) 是从 E 到 F 的部分算子， G 是 E 的向量子空间，则 \(u\) 到 G 上的约化是从 E 到 F 的部分算子（ \(\Gamma \cap (\mathrm{G} \times \mathrm{F}), \mathrm{E}, \mathrm{F}\) ）。它的定义域是 \(\operatorname{dom}(u) \cap \mathrm{G}\) ；当不致与 \(u|_{\mathrm{G}}\) 到子空间 G 上的限制相混淆时，有时把它记作 \(u\) 。
\item
  设 \(v\) 是从 F 到 E 的单射部分算子。这时 \(v^{-1} = (\Gamma^{-1}, \mathrm{E}, \mathrm{F})\) 的逆对应 \(v\) 是一个部分算子，满足 \(\operatorname{dom}(v^{-1}) = \operatorname{Im}(v)\) 。称 \(v^{-1}\) 为 \(v\) 的逆部分算子。我们有等式 \(v \circ v^{-1} = 1_{\operatorname{dom}(v^{-1})}\) （在 \(\mathcal{P}(\mathrm{E})\) 中）与 \(v^{-1} \circ v = 1_{\operatorname{dom}(v)}\) （在 \(\mathcal{P}(\mathrm{F})\) 中）。部分算子 \(v^{-1}\) 是单射的，且有 \((v^{-1})^{-1} = v\) 。
\item
  设 E 、 F 与 G 是向量空间。设 u （相应地， v ）是从 E 到 F （相应地，从 F 到 G ）的单射部分算子。那么部分算子 \(v \circ u\) 是单射的，且 \((v \circ u)^{-1} = u^{-1} \circ v^{-1}\) 。
\item
  若 u 与 v 是从 E 到 F 的部分算子，我们说 v 是 u 的扩张，并记 \(u \subset v\) ，如果 u 的图像包含在 v 的图像中。这蕴含 \(\text{dom}(u) \subset \text{dom}(v)\) ，且 u 是 v 到 \(\text{dom}(u)\) 上的约化。关系 “ \(u \subset v\) ” 是 \(\mathcal{P}(\text{E};\text{F})\) 上的一个序关系。例如，对所有 \(au \subset ua\) 与每个 \(a \in K\) ，有 \(u \in \mathcal{P}(\text{E};\text{F})\) 。
\item
  设 E 是 K 上的向量空间， \((\mathrm{F}_{i})_{i\in\mathrm{I}}\) 是 K 上向量空间的一个族。对 \(i\in I\) ，设 \(u_{i}\) 是从 E 到 \(F_{i}\) 的一个部分算子。诸 \(u_{i}\) 的乘积部分算子是从 E 到诸空间 \(F_{i}\) 的乘积向量空间中的部分算子，其定义域是诸 \(\mathrm{dom}(u_{i})\) 的交 D ，它把 \(x\in D\) 对应到族 \((u_{i}(x))_{i\in\mathrm{I}}\) 。把它记作 \((u_{i})_{i\in\mathrm{I}}\) 。
\item
  设 A: F×F → F 是线性映射 \((x,y)\mapsto x+y\) 。设 u 与 v 是从 E 到 F 的部分算子。和 \(u+v\) 是从 E 到 F 的部分算子 A ∘ \((u,v)\) 。它的定义域是 dom \((u)\cap\text{dom}(v)\) 。对 \(\mathcal{P}(\text{E};\text{F})\) 中的 u 、 v 、 w ，有 \((u+v)+w=u+(v+w)\) 。
\end{enumerate}

设 G 是 K 上的向量空间。对 \(\mathcal{P}(\mathrm{E};\mathrm{F})\) 中所有 u 与 v 以及每个 \(w \in \mathcal{P}(\mathrm{F};\mathrm{G})\) ，有 \(w \circ u + w \circ v \subset w \circ (u + v)\) 。一般而言，这一公式中没有等号（IV, p. 344 的习题 1），但当 w 的定义域等于 F 时就是这种情形。对 \(w \in \mathcal{P}(\mathrm{G};\mathrm{E})\) ，有 \(u \circ w + v \circ w = (u + v) \circ w\) 。

\begin{enumerate}
\def\labelenumi{(\roman{enumi})}
\setcounter{enumi}{6}
\tightlist
\item
  设 L 是域 K 的一个扩张。设 E 与 F 是 K 上的向量空间， \(\mathrm{E}_{(\mathrm{L})} = \mathrm{L} \otimes_{\mathrm{K}} \mathrm{E}\) ， \(\mathrm{F}_{(\mathrm{L})} = \mathrm{L} \otimes_{\mathrm{K}} \mathrm{F}\) 是由 K 到 L 的纯量扩张得到的 L-向量空间（A, II, p. 82）。对每个从 E 到 F 的部分算子 u ，我们用 \(u_{(L)}\) 表示从 \(\mathrm{E}_{(\mathrm{L})}\) 到 \(\mathrm{F}_{(\mathrm{L})}\) 中的、其图像为 \(L \otimes_{K} \Gamma_{u}\) （ \(\mathrm{E}_{(\mathrm{L})} \times \mathrm{F}_{(\mathrm{L})}\) 的向量子空间）的部分算子；它的定义域是 \(L \otimes_{K} \operatorname{dom}(u)\) ，且它在这一定义域上与把 \(1 \otimes x\) 映到 \(1 \otimes u(x)\) （对所有 \(x \in \operatorname{dom}(u)\) ）的唯一线性映射重合。
\end{enumerate}

设 v 是从 E 到 F 的部分算子。我们有 \(u \subset v\) 当且仅当 \(u_{(L)} \subset v_{(L)}\) 。

\begin{enumerate}
\def\labelenumi{(\roman{enumi})}
\setcounter{enumi}{7}
\tightlist
\item
  设 \(E_{1}\) ， \(F_{1}\) ， \(E_{2}\) ， \(F_{2}\) 都是 K-向量空间。设 u （相应地， v ）是从 \(E_{1}\) 到 \(F_{1}\) （相应地，从 \(E_{2}\) 到 \(F_{2}\) ）的部分算子。我们用 \(u \otimes v\) 表示从 \(E_{1} \otimes F_{1}\) 到 \(E_{2} \otimes F_{2}\) 中的、定义域为 dom(u) \(\otimes\) dom(v) 且满足 \((u \otimes v)(x \otimes y) = u(x) \otimes v(y)\) （对所有 \((x, y) \in E_{1} \times E_{2}\) ）的部分算子。
\end{enumerate}

